% Generated draft: requires mathematical and source review.
\documentclass[11pt]{article}
\usepackage[margin=25mm]{geometry}
\usepackage{fontspec}
\setmainfont{Latin Modern Roman}
\usepackage{amsmath,amssymb,mathtools,unicode-math}
\setmathfont{Latin Modern Math}
\usepackage{enumitem,needspace,xurl,hyperref}
\setlength{\parindent}{0pt}
\setlength{\parskip}{5pt}
\setlength{\emergencystretch}{3em}
\newcommand{\N}{\mathbb{N}}
\newcommand{\Z}{\mathbb{Z}}
\newcommand{\Q}{\mathbb{Q}}
\newcommand{\R}{\mathbb{R}}
\newcommand{\C}{\mathbb{C}}
\newcommand{\F}{\mathbb{F}}
\newcommand{\A}{\mathbb{A}}
\newcommand{\PP}{\mathbb{P}}
\newcommand{\E}{\mathbb{E}}
\newcommand{\eps}{\varepsilon}
\newcommand{\dd}{\,\mathrm{d}}
\begin{document}
\section*{UnsolvedMath ID: 1101219}
\textbf{Source problem code:} catalogue entry
\subsection{Definitions and mathematical statement}
Adopt the domain of the supplied research summary: \(S_g\) is a closed connected orientable surface of genus \(g\ge2\), and \(G=\operatorname{MCG}(S_g)\) has a fixed word metric \(d\). A quasi-isometry \(\varphi:G\to G\) satisfies, for some \(L\ge1,C,D\ge0\),
\[
L^{-1}d(x,y)-C\le d(\varphi(x),\varphi(y))\le Ld(x,y)+C,
\qquad G\subseteq N_D(\varphi(G)).
\]
Here \(N_R(Y)=\{x:d(x,Y)\le R\}\). A maximal flat means the group-theoretic flat associated with a free abelian subgroup of maximal possible rank: concretely a left coset \(hA\), with \(A\le G\) free abelian of that rank, considered up to finite Hausdorff distance. For subsets \(U,V\), \(d_{\mathrm H}(U,V)\le R\) means \(U\subseteq N_R(V)\) and \(V\subseteq N_R(U)\).

The coarse flat-preservation question is whether every such \(\varphi\) has the following property:
\[
\forall\text{ maximal flats }F\subseteq G\quad
\exists\text{ maximal flat }F'\subseteq G:\quad
 d_{\mathrm H}(\varphi(F),F')<\infty.
\]
The supplied research summary records an affirmative answer, with a uniform bound
\[
\exists R=R(L,C,D)\ \forall F\ \exists F':
 d_{\mathrm H}(\varphi(F),F')\le R,
\]
for the fixed surface and word metric. This is preservation up to bounded Hausdorff distance, the meaning supplied by the summary, rather than literal equality of subsets. The summary attributes it to BKMM Theorem 10.3 combined with the Birman--Lubotzky--McCarthy classification of maximal-rank abelian subgroups, and reports an independent proof implicit in Hamenst\"adt's preprint.
\subsection{Short English statement}
Does every quasi-isometry of a closed surface mapping class group carry every maximal-rank abelian flat within a uniformly bounded Hausdorff distance of another such flat? The supplied summary answers yes for genus at least two.
\subsection{Sources}
\begin{itemize}
\item Local UnsolvedMath registry, permanent ID \(1101219\), exact source code \texttt{catalogue entry}, ``Questions in Geometric Group Theory --- Q 12.19.''
\item Bestvina, \emph{Questions in Geometric Group Theory} (2004), Question 12.19, PDF page 23. \url{https://www.math.utah.edu/~bestvina/eprints/questions-updated.pdf}; manual transcription; recorded access 2026-07-29.
\item Supplied research summary: BKMM, Theorem 10.3, \emph{Geometry \& Topology} \textbf{16} (2012), together with Birman--Lubotzky--McCarthy's classification; Hamenst\"adt, quasi-isometric rigidity preprint, arXiv:math/0512429. Further bibliographic details are not supplied.
\end{itemize}
\end{document}
